%! Vector Spaces and Subspaces — Unit 3, Lesson 3.1 — Warm-Up (Answer Key)
\documentclass[10pt]{article}
\usepackage{linalg-article}
\usepackage{linalg-key}   % pulls in -boxes; do NOT also load -boxes
\newcommand{\TallMath}[1]{$\displaystyle #1\rule[-1.4em]{0pt}{3.2em}$}
\newcommand{\xx}{\mathbf{x}}
\newcommand{\yy}{\mathbf{y}}
\newcommand{\zero}{\mathbf{0}}

\begin{document}

\pageheader{Unit 3, Lesson 3.1}{Warm-Up --- Answer Key}
\namedateperiod

\vspace{0.12in}

\begin{notesbox}{Getting Ready: Skills We'll Use Today}
\small These skills from Units 1 and 3.0 power today's proof. Answer quickly.

\vspace{4pt}
\textbf{1. $A\xx$ is a combination of the columns.} With $A=\begin{bmatrix} 1 & 2 \\ 3 & 1 \end{bmatrix}$
and $\xx=\begin{bmatrix} 2 \\ 1 \end{bmatrix}$, compute $A\xx = 2\begin{bsmallmatrix} 1\\3 \end{bsmallmatrix}+1\begin{bsmallmatrix} 2\\1 \end{bsmallmatrix}$:
\[
  A\xx = \begin{bmatrix}{\color{keyred}\mathbf{4}}\\[2pt]{\color{keyred}\mathbf{7}}\end{bmatrix}.
\]

\vspace{4pt}
\textbf{2. The distributive law $A(\xx+\yy)=A\xx+A\yy$.} Keep the same $A$, and take
$\yy=\begin{bmatrix} 1 \\ 0 \end{bmatrix}$. Fill both sides and check they match.
\[
  A(\xx+\yy)=A\begin{bsmallmatrix} 3\\1 \end{bsmallmatrix}=\begin{bmatrix}{\color{keyred}\mathbf{5}}\\[2pt]{\color{keyred}\mathbf{10}}\end{bmatrix},
  \qquad
  A\xx+A\yy=\begin{bsmallmatrix} 4\\7 \end{bsmallmatrix}+\begin{bmatrix}{\color{keyred}\mathbf{1}}\\[2pt]{\color{keyred}\mathbf{3}}\end{bmatrix}
  =\begin{bmatrix}{\color{keyred}\mathbf{5}}\\[2pt]{\color{keyred}\mathbf{10}}\end{bmatrix}.
\]
Are the two sides equal? \; \ans{yes} \quad / \quad no

\vspace{4pt}
\textbf{3. Is $\zero$ in the set?} (A fast reject from Lesson~3.0.) Is the zero vector
$\begin{bsmallmatrix} 0\\0 \end{bsmallmatrix}$ on the line $y=2x-1$? \ansline{No --- at $x=0$ the line gives $y=-1$, not $0$. So the set fails the fast $\zero$ check and cannot be a subspace.}
\end{notesbox}

\begin{teachernote}
Item~2 is the engine of today's lesson: $A(\xx+\yy)=A\xx+A\yy$ is exactly why adding two reachable
targets gives a reachable target, so the column space is closed. Item~3 is the quick reject we
formalize as the subspace requirement.
\end{teachernote}

\end{document}
